\documentclass[11pt]{article}

\usepackage[margin=1in]{geometry}
\usepackage{amsmath}
\usepackage{booktabs}
\usepackage{graphicx}
\usepackage{natbib}
\usepackage[hidelinks]{hyperref}

\title{Educational Upgrading and Wage Growth: Evidence across Countries}
\author{Oliver Pardo}
\date{\today}

\begin{document}

\maketitle

\begin{abstract}
This paper separates changes in average employee wages into two exactly
additive components: changes in educational attainment and changes in wages
within educational groups. We combine annual ILOSTAT series on average monthly
earnings with employee counts for four harmonized education groups. A
preliminary source-matched audit identifies 117 countries with at least one
complete observation and 44 countries with at least five observations from the
same statistical source that pass the initial quality screen. The paper will
document how the two components vary across countries and over time, without
interpreting either component as a causal return to education or as a direct
measure of productivity.
\end{abstract}

\section{Introduction}

Average wages can rise because workers attain higher levels of education or
because employees with a given educational attainment receive higher wages.
These two changes have different interpretations. Educational upgrading changes
the composition of employment, whereas wage growth within educational groups
captures changes that occur after holding that composition fixed.

This paper measures the contribution of each component to wage growth across
countries. It asks three questions. First, how much of the change in average
real wages is associated with educational upgrading? Second, how much is
associated with changes in wages within educational groups? Third, how do these
contributions vary with countries' initial education and income levels and
across major episodes, including the pandemic and the subsequent recovery?

The analysis uses annual ILOSTAT data for employees. This restriction is
substantive rather than terminological: earnings refer to employees, so the
population weights must also refer to employees. We therefore match average
monthly earnings by education to employment by education and status in
employment, retaining the employee category and the same country, year, and
statistical source.

\section{Related literature}

The accounting framework follows the symmetric decomposition proposed by
\citet{kitagawa1955}. It is related to the broader decomposition literature
reviewed by \citet{fortinlemieuxfirpo2011}, but differs from regression-based
Oaxaca--Blinder exercises because it decomposes the observed mean directly.
The paper also relates to evidence on compositional effects in aggregate wage
growth \citep{dalyhobijn2017,kouvavasetal2019} and to studies of educational
upgrading and skill premia.

\section{Data}

\subsection{ILOSTAT earnings and employment}

The remuneration measure is average nominal monthly earnings of employees by
sex and education. Population weights come from employment by sex, status in
employment, and education. The baseline sample retains total sex, employees,
and four aggregate education groups: less than basic, basic, intermediate, and
advanced.

The two tables are matched by country, year, and source. The initial audit finds
no duplicate observations at that grain. It also compares the weighted average
of the four education groups with ILOSTAT's published total and records the
share of employees whose educational attainment is not stated. Country-years
with a reconstruction discrepancy above 5 percent or an unknown-education
share above 5 percent are excluded from the preliminary clean sample.

\subsection{Prices and sample construction}

Nominal earnings will be converted into constant local-currency units using a
country-specific consumer price index. The decomposition is conducted within
countries, so purchasing-power-parity conversion is not required for the
baseline growth estimates. Cross-country figures will report contributions as
growth rates or shares of the total change.

\section{Method}

\subsection{Exact decomposition}

Let \(R_t\) denote average remuneration in period \(t\). The population and
remuneration concept are defined separately in each paper. Educational
attainment partitions that population into mutually exclusive groups indexed by
\(g\):
\begin{equation}
  R_t = \sum_g s_{g,t} r_{g,t},
  \label{eq:remuneration_identity}
\end{equation}
where \(s_{g,t}\) is group \(g\)'s share of the relevant population and
\(r_{g,t}\) is its average remuneration.

For two periods, the change in average remuneration can be written as
\begin{equation}
  R_1-R_0 =
  \sum_g \bar{s}_g\left(r_{g,1}-r_{g,0}\right)
  +
  \sum_g \bar{r}_g\left(s_{g,1}-s_{g,0}\right),
  \label{eq:symmetric_decomposition}
\end{equation}
where
\(\bar{s}_g=(s_{g,1}+s_{g,0})/2\) and
\(\bar{r}_g=(r_{g,1}+r_{g,0})/2\).
The first term measures changes in remuneration within educational groups. The
second measures changes in the educational composition of the relevant
population. This symmetric decomposition is exactly additive and does not
depend on choosing either period as the reference.

Equation~\eqref{eq:symmetric_decomposition} is an accounting identity. The
within-group component is not a measure of productivity, and the composition
component does not identify the causal effect of education. Changes in either
component may reflect selection, labor supply, labor demand, institutions,
hours, occupations, sectors, or survey coverage.


\subsection{Comparison windows}

The baseline specification will use source-consistent country panels with at
least five valid annual observations. Results will be reported for common
windows when coverage permits and for each country's longest valid interval as
a robustness exercise. Comparisons will not cross documented breaks in series
without an explicit sensitivity test.

\section{Results}

\subsection{Educational upgrading}

\subsection{Within-education wage growth}

\subsection{Cross-country heterogeneity}

\subsection{The pandemic and recovery}

\section{Robustness}

\section{Conclusion}

\bibliographystyle{apalike}
\bibliography{references}

\end{document}
